% !PRINTED 659
\weberpage{659}

Es kommt nun f\"ur den Beweis des Haupttheorems I. alles auf den
Beweis der beiden S\"atze an, die aus den Eigenschaften der
Functionale $\varphi$ und $\varepsilon$ zu folgern sind:

A. \gesperrt{Die Functionale $\varphi_n$ geh\"oren der Hauptclasse an.}

K\"onnen wir dieses Lemma voraussetzen, so ist $\varphi_n$ mit
einer Zahl $\alpha_n$ associirt, und wir k\"onnen die Formel (14)
mit etwas ver\"anderter Bedeutung von $\varepsilon_n$ so
darstellen:
\weq{15}{\omega_n = \varepsilon_n \alpha_n^m \prod^{p} (r^{m_2 n}, \eta)^m.}
In dieser Formel ist aber $\varepsilon_n$ eine \gesperrt{numerische Einheit}
des K\"orpers $\Omega_m$, die dem Satze 9. gen\"ugt.

K\"onnen wir also aus dem Satze 9. noch das zweite Lemma
beweisen:

B. \gesperrt{Eine numerische Einheit $\varepsilon_n$ des K\"orpers
$\Omega_m$, die dem Satze 9. gen\"ugt, ist die $m^{\text{te}}$
Potenz einer Einheit in $\Omega_m$, multiplicirt mit einer Potenz
von $r$;}

so k\"onnen wir in (15) die $m^{\text{te}}$ Wurzel ziehen, und
erhalten, wenn mit $\varrho$ eine Einheitswurzel von
m\"oglicherweise h\"oherem Grade als $m$ und mit $\alpha$ eine
Zahl des K\"orpers $\Omega_m$ bezeichnet wird:
\weq{16}{(r, x_0) = \varrho \alpha \prod^{p} (r^{m_2}, \eta),}
wo nun rechts lauter Kreistheilungszahlen stehen und wodurch
daher das Theorem I. erwiesen ist.

\begin{center}
\S.\,182.

\gesperrt{Beweis des ersten H\"ulfssatzes f\"ur ein ungerades $m$.}
\end{center}

Wir haben im vorigen Paragraphen den Beweis des Theorems I. von
zwei H\"ulfss\"atzen abh\"angig gemacht, deren Beweis nun ferner
zu suchen ist.

Wir formuliren den ersten dieser H\"ulfss\"atze so:

a) \gesperrt{Ist $\varphi_n$ ein System von Null verschiedener
conjugirter Functionale des K\"orpers $\Omega_m$, von denen
bekannt ist, dass}
\[
\varphi_n^m \text{ \gesperrt{und} } \varphi_n^{-1} \varphi_1^n
\]
\gesperrt{der Hauptclasse angeh\"oren, so geh\"oren die
Functionale $\varphi$ selbst der Hauptclasse an.}
